\chapter{个体学习}\label{s:individual}

前面几章探讨了老师可以做些什么来帮助学习者。
这一章看的则是学习者自己可以做些什么，
办法是改变学习策略，并保证充足的休息。

最有效的策略，是从\gref{g:passive-learning}{被动学习}
转向\gref{g:active-learning}{主动学习}~\cite{Hpl2018}，
这能显著提升表现、降低不及格率~\cite{Free2014}：

\begin{longtable}{ll}
\textbf{被动}		& \textbf{主动} \\
阅读某个主题		& 做练习 \\
看视频			& 讨论一个话题 \\
听讲座			& 试着把它讲出来
\end{longtable}

\noindent
回到我们那个简化的认知架构模型（\figref{f:arch-model}），\index{认知架构}
主动学习之所以更有效，是因为它让新信息在短期记忆里停留得更久，\index{短期记忆}
从而增加了它被成功编码并存入长期记忆的可能性。\index{长期记忆}
而且，通过在新信息到来时就使用它，
学习者在这些信息和他们已知的东西之间建立或强化了联系，
这反过来又增加了他们日后能够回忆起它的几率。

把学习收获最大化的另一个关键，是\gref{g:metacognition}{元认知}，
也就是对自己的思考进行思考。
正如好音乐家会听自己的演奏，
好老师会反思自己的教学（\chapref{s:performance}），
学习者如果会做计划、
设定目标、
并监控自己的进展，也会学得更好、更快。
学习者很难抽象地掌握这些技能——光告诉他们「要做计划」是没有任何效果的——
但课程可以被设计成鼓励良好的学习习惯，
而在课堂上点出这些做法，
能帮学习者意识到学习是一项像其他技能一样可以提升的技能~\cite{McGu2015,Miya2018}。

最大的奖项是\gref{g:transfer-of-learning}{学习迁移}，
也就是我们学到的某样东西帮助我们更快地学会其他东西。
研究者区分了\grefdex{g:near-transfer}{近迁移}{学习迁移!近迁移}，
它发生在相似或相关的领域之间，比如数学里的分数和小数；
以及\grefdex{g:far-transfer}{远迁移}{学习迁移!远迁移}，
它发生在不相干的领域之间——例如，
学下棋有助于数学推理（或反过来）这种说法。

近迁移无疑会发生——如果没有它，
除了简单的记忆之外，就什么学习都不可能发生——
老师也一直在利用它，
办法是给学习者做一些与课上刚讲过的内容相似的练习。
然而，
\cite{Sala2017}分析了许多关于远迁移的研究，
得出结论：

\begin{quote}

  {\ldots}结果显示效应为小到中等。
  不过，效应量与实验设计的质量成反比{\ldots}
  我们得出的结论是，学习的远迁移很少发生。

\end{quote}

当远迁移确实发生时，
它似乎只有在某个学科被掌握之后才会出现~\cite{Gick1987}。
在实践中，
这意味着学编程不会帮你下棋，反之亦然。

\seclbl{六种策略}{s:individual-strategies}

心理学家用各种各样的方法研究学习，
但就什么真正有效，得出了相似的结论~\cite{Mark2018}。
\hreffoot{http://www.learningscientists.org/}{Learning Scientists}
把这些策略中的六种整理了出来，
并总结在\hreffoot{http://www.learningscientists.org/downloadable-materials}{一套可下载的海报}里。
把这些策略教给学习者，
并在课堂上用到它们时点出它们的名字，
能帮他们学会如何学得更快、更好~\cite{Wein2018a,Wein2018b}。

\subsection*{间隔练习}
\index{间隔练习（学习策略）}
\index{学习策略!间隔练习}

把十小时的学习分散在五天里，
比两个五小时的日子更有效，
也比一个十小时的日子好得多。
因此你应该制定一个把学习活动分散到一段时间里的学习日程：
每天至少留出半小时来学每个主题，
而不是在考前一晚把所有东西囫囵塞进去~\cite{Kang2016}。

你还应该在每次课后复习材料，
但不要马上复习——至少休息半小时。
复习时，
一定要至少带上一点点更早的材料：
例如，
花二十分钟看当天课上的笔记，
然后分别花五分钟看前一天和一周前的材料。
这样做还能帮你在还来得及改正或提问的时候，
发现之前几套笔记里的空缺或错误：
考前一晚才发现自己完全不知道为什么当初把「解调！！」划了三道线，是很痛苦的。

复习时，
把那些你忘掉的东西记下来：
例如，为每个你记不住、
或者记错了的事实做一张卡片~\cite{Matt2019}。
这能帮你把下一轮学习的重点放在最需要关注的地方。

\begin{aside}{讲座的价值}
  根据~\cite{Mill2016a}，
  「面对面课程里占主导地位的讲座，是相对低效的教学方式，
  但它们或许有助于把材料分散到一段时间里，
  因为讲座会按固定的日程随时间展开。
  相比之下，
  取决于课程是怎么设置的，
  在线学生有时可以拖到作业临近之前，才接触到材料。」
\end{aside}

\subsection*{提取练习}
\index{提取练习（学习策略）}
\index{学习策略!提取练习}

长期记忆的限制因素不是保持（存了什么），
而是提取（能取出什么）。
对具体信息的提取会随练习而改善，
所以在真实情境中的表现，可以通过做练习测验、
或凭记忆总结某个主题的细节、然后再核对记住了什么没记住什么来提升。
例如，
\cite{Karp2008}发现，反复测验把词表的回忆率从 35\% 提升到了 80\%。

当练习使用的活动与测验相似时，提取效果更好。
例如，
写个人日记有助于多项选择题测验，
但不如做练习测验有效~\cite{Mill2016a}。
这一现象叫做
\gref{g:transfer-appropriate-processing}{迁移适宜加工}。

锻炼提取技能的一种办法，是把问题解两遍。
第一遍，
完全凭记忆来做，不看笔记，也不和同伴讨论。
对照老师给的评分标准给自己打分之后，
再用你想用的任何资源把问题解一遍。
两遍之间的差距，就显示了你提取并运用知识的能力有多好。

另一种办法（上面提到过）是做卡片。
纸质卡片一面写问题或其他提示，另一面写答案，
手机上也有很多卡片应用。
如果你是结伴或结组学习，
和同伴交换卡片，
能帮你发现一些你可能漏掉或误解了的重要想法。

\grefdex{g:read-cover-retrieve}{阅读—遮盖—回忆}{阅读—遮盖—回忆}
是卡片的一种快捷替代。
你在看某样材料时，
用小便利贴把关键词或段落盖住。
看完之后，
再从头过一遍，看看你能猜出每张便利贴下面是什么猜得多准。
不管你用哪种方法，
都不要只练习回忆事实和定义：
一定要同时检查自己对大概念、
以及它们之间联系的理解。
画一张概念图，然后把它和你的笔记、
或之前画过的概念图作比较，
是快速做到这一点的办法。

\begin{aside}{超校正}
  学习研究中一个很有力的发现是
  \gref{g:hypercorrection}{超校正效应}~\cite{Metc2016}。
  大多数人都不喜欢被告知自己错了，
  所以人们有理由假设：
  一个人对他在测验中给出的答案越有信心，
  如果他其实错了，就越难改变他的想法。
  结果事实恰恰相反：
  一个人越是确信自己是对的，
  一旦被纠正，就越有可能不再重复这个错误。
\end{aside}

\subsection*{交错练习}
\index{交错练习（学习策略）}
\index{学习策略!交错练习}

给练习留出间隔的一种办法，是把不同主题的学习交错起来：
不要先把一个科目学透，再学第二个、第三个，
而是把学习时段打乱。
更进一步，
把顺序也换一换：
A-B-C-B-A-C 比 A-B-C-A-B-C 好，
而后者又比 A-A-B-B-C-C 好~\cite{Rohr2015}。
这样做之所以有效，是因为交错有助于在不同主题之间建立更多联系，
进而改善提取。

每个条目该花多长时间，
取决于科目以及你对它有多熟。
10 到 30 分钟之间的某个长度，足够让你进入
心流状态（\secref{s:individual-time}），
但又不至于让你走神。
交错学习一开始会比一次只专注一个主题感觉更难，
但那是它正在起作用的信号。
如果你用卡片或练习测验来衡量进展，
只需几天就应该看到进步。

\subsection*{精细加工}
\index{精细加工（学习策略）}
\index{学习策略!精细加工}

在过材料的过程中向自己解释，
有助于你理解和记住它们。
一种做法是，在练习测验里每个答案之后，
补上一段说明，解释为什么这个答案是对的，
或者反过来解释为什么另一个看似合理的答案不对。
另一种做法是，向自己说明一个新想法
与你之前见过的某个想法有什么相似或不同。

自言自语看起来像是种奇怪的学习方式，
但~\cite{Biel1995}发现，
受过自我解释训练的人比没有受过训练的人表现更好。
同样，
\cite{Chi1989}发现，有些学习者在解决问题时，
一旦碰到看不懂的步骤就干脆停住了。
另一些人则会暂停下来，对正在发生的事情作出解释，
后一组人学得更快。
锻炼这项技能的一个练习，是和一个班一起逐行过一遍示例程序，
让不同的人轮流解释每一行，
并说明它为什么在那里、它达成了什么。

\subsection*{具体例子}
\index{具体例子（学习策略）}
\index{学习策略!具体例子}

精细加工的一种特别有用的形式，是使用具体例子。
每当你有一条一般原理的陈述，
就试着给出一个或多个它应用的例子，
或者反过来，拿每个具体问题，列出它所体现的一般原理。
\cite{Raws2014}发现，像这样把例子和定义交错起来，
更可能让学习者正确记住定义。

一种结构化的做法是
\hreffoot{https://betterexplained.com/articles/adept-method/}{ADEPT 方法}：
\index{ADEPT（课程模式）}
\index{课程模式!ADEPT}
给出一个\textbf{类比}（Analogy）、
画一张\textbf{图}（Diagram）、
举一个\textbf{例子}（Example）、
用\textbf{平实}的语言（Plain language）描述这个想法，
然后给出\textbf{技术}细节（Technical details）。
再说一遍，
如果你是和一个同伴或一个小组一起学习，
你们可以交换并互相检查：
看看你们是否都同意别人的例子确实体现了所讨论的原理，
或者他们没列出来的例子里用到了哪些原理。

另一种有用的技巧是对比教学，
\index{对比教学（课程模式）}
\index{课程模式!对比教学}
也就是向学习者展示一个解\emph{不}是什么，
或者某个技巧\emph{解决不了}哪类问题。
例如，
教孩子化简分数时，
一定要给他们几个像 5/7 这样无法化简的分数，
好让他们不会因为寻找根本不存在的答案而受挫。

\subsection*{双重编码}
\index{双重编码（学习策略）}
\index{学习策略!双重编码}

\hreffoot{http://www.learningscientists.org/}{Learning Scientists} 描述的
六种核心策略中的最后一种，
是把文字和图像一起呈现。
正如\secref{s:architecture-brain}里讨论的，
我们大脑中不同的子系统处理和存储语言信息与视觉信息，
所以如果通过两条通道呈现互补的信息，
它们就能相互强化。
然而，
当同样的信息同时在两条不同的通道里呈现时，学习反而更差，
因为这时大脑得花力气去核对两条通道是否一致~\cite{Maye2003}。

利用双重编码的一种办法，是画或标注时间线、
地图、家谱，
或者其他对材料来说似乎合适的东西。
（我个人很喜欢那种显示程序里哪个函数调用了哪个函数的图。）
先画一张\emph{没有}标注的图，
过一阵再回来给它加标注，
是极好的提取练习。

\seclbl{时间管理}{s:individual-time}

我曾经吹嘘自己工作了多少小时。
当然不是明说——我\emph{好歹}有点社交技巧——
但我会在中午前后不修边幅、打着哈欠来上课，
随口跟任何愿意听的人提起
我一直干到凌晨六点。

回头想想，
我想不起自己当时是想让谁刮目相看。
我记住的反而是，
那些通宵干的活里，有多少在我睡了一觉之后被我扔掉了，
而那些没扔掉的，又给我的成绩造成了多大的损害。

我的错误在于把「工作」和「有产出」混为一谈。
你不下点功夫就生产不出软件（或任何别的东西），
但你可以轻易地干很多活、却产不出任何有价值的东西。
要让人相信这一点可能很难，
尤其当他们只有一二十岁的时候，
但这样做回报巨大。

对过度劳累和睡眠剥夺的科学研究至少可以追溯到 1890 年代——参见
\cite{Robi2005}这篇简短易读的综述。\index{过度劳累}\index{睡眠剥夺}
对学习者来说最重要的结果有：

\begin{enumerate}

\item
  一天工作超过 8 小时并持续一段时间，
  会降低你的总产出，
  而不只是每小时的产出——也就是进入连轴转状态时，你完成的总量（不只是每小时）反而更少。

\item
  一口气工作超过 21 小时，
  你犯下灾难性错误的几率，和达到法定醉酒标准时一样高。

\item
  产出率在一天之中会变化，
  最高的产出出现在最初的 4 到 6 小时。
  过了足够多的钟点之后，
  产出率趋近于零；
  最后会变成负数。

\end{enumerate}

这些事实已经被重复验证了一个多世纪，
它们背后的数据和把吸烟与肺癌联系起来的那些数据一样扎实。
问题在于，
\emph{人们通常察觉不到自己的能力在下降}。
就像醉汉以为自己还能开车一样，
被剥夺睡眠的人意识不到
自己话都说不完整了（更别提想不完整）。
每周五个 8 小时的工作日，已被证明能在所有被研究过的行业中
最大化长期总产出；
学习或编程也不例外。

那偶尔来一下短时的冲刺呢，
比如为了赶截止日期熬个通宵？
这也被研究过，
结果并不好看。
你每清醒 24 小时，思考能力就下降 25\%。
换种说法，
熬一个通宵之后，普通人的智商只有 75，
这把他们放到了人口中垫底的 5\%。
如果你连熬两个通宵，你的有效智商就是 50，
到那个水平，人们通常会被判定为无法独立生活。

「可是——可是——我有那么多作业要做！」你说。
「而且它们全都同时到期！
我\emph{必须}加班加点才能全做完！」
不对：
人为了有产出，必须分清轻重缓急、保持专注，
而为了做到这一点，
他们得先学会怎么做。
一种被广泛使用的技巧，是把要做的事情列个清单，
按优先级排序，
然后关掉邮件和其他干扰 30—60 分钟，
完成其中一项任务。
如果待办清单里任何一项任务超过一小时，
就把它拆成更小的块，再分别排优先级。

这件事最重要的部分，是关掉干扰。
不管很多人多么愿意相信，
人类并不擅长多任务。\index{多任务}
我们能变得擅长的，是\gref{g:automaticity}{自动化}，
也就是在做一件事的同时，能在后台做另一件例行公事的能力~\cite{Mill2016a}。
我们大多数人都能一边切洋葱一边说话，
或者一边喝咖啡一边看书；
经过练习，
我们也能一边听一边记笔记，
但我们没法在关注别的东西的同时有效地学习、
编程、
或做其他费脑子的任务——我们只是以为自己能。

组织和准备的意义，
是为了进入尽可能高效的心理状态。
心理学家把它叫做\gref{g:flow}{心流}~\cite{Csik2008}；
运动员叫它「进入状态」，
音乐家则说他们沉浸在演奏之中。
不管你怎么称呼它，
人们在这种状态下单位时间的产出远高于平常。
坏消息是，
无论干扰多短暂，一次干扰之后大约都需要 10 分钟才能回到心流状态。
这意味着，如果你每小时被打断六次，
你就\emph{永远}达不到产出的巅峰。

\newpage

\begin{aside}{他怎么知道的？}

  在他 1961 年的短篇小说「\hreffoot{https://en.wikipedia.org/wiki/Harrison\_Bergeron}{哈里森·伯杰龙}」里，
  库尔特·冯内古特描述了一个人人都被强制平等的未来。
  长得好看的人必须戴面具，
  体格强健的人必须负重——而聪明人
  则被强制携带会随机打断他们思绪的收音机。
  我有时会想，是不是——哦，等等，我的手机刚——抱歉，我们刚才在说什么？

\end{aside}

\seclbl{同伴评价}{s:individual-peer}
\index{同伴评价}

让团队里的人给同伴打分，是工业界的一种常见做法。
\cite{Sond2012}综述了关于同伴评价的文献，
区分了打分与评阅。
他们发现，同伴评价增加了反馈的数量、多样性和及时性，
帮助学习者练习更高层次的思考，
鼓励反思性实践，
并支持社交技能的发展。
其中的担忧也不难预料：
效度和信度、
动机和拖延、
捣乱者、串通和抄袭。

然而，
证据表明，这些担忧在大多数课堂里并不显著。
例如，
\cite{Kauf2000}在两门本科工程课程中，
从几个维度比较了保密的同伴评分与分数，
发现自评和同伴评分在统计上一致、
串通并不显著（也就是人们并没给所有同伴都打高分）、
学习者没有给自己虚高评分，
而且关键在于，
评分没有受到性别或种族的偏见影响。

落实同伴评价的一种方式，是\gref{g:contributing-student}{贡献型学生教学法}，
学习者制作作品，为其他人的学习作出贡献。
这可以是开发一节短课并与全班分享、
向题库里添加题目、
或者为课上发布而撰写某次讲座的笔记。
例如，
\cite{Fran2018}发现，制作短视频给同伴讲概念的学习者，
自己的学习相比只学材料或看视频的人有显著提升。
我发现，让学习者每天和全班分享一个 bug 及其修复，
有助于提升他们的分析能力，并减轻冒充者综合征。

另一种做法是\gref{g:calibrated-peer-review}{校准式同伴互评}，
学习者用评分标准评阅一个或多个样例，
并把自己的评价与老师对同一份作品的评阅作比较~\cite{Kulk2013}。
一旦学习者的评价与老师的足够接近，
他们就开始评阅同伴的真作。
如果把几个同伴的评价综合起来，
其准确性可以和老师评阅相媲美~\cite{Pare2008}。

和别的一切一样，
评价也离不开评分标准。
\secref{s:checklists-teameval}里的评价表给了一个样例，供你上手。
使用它时，
给你自己和每位队友排名，
然后计算并比较分数。
差距很大通常说明需要更长时间的一次谈话。

\seclbl{练习}{s:individual-exercises}

\exercise{学习策略}{个人}{20}

\begin{enumerate}

\item
  六种学习策略中，你经常用哪几种？
  哪几种不用？

\item
  写下你希望学习者掌握的三个一般概念，
  并为每个概念给出两个具体例子
  （练习具体例子）。
  对每个概念，
  从你给的某个例子往回推，解释这个概念是如何说明它的
  （精细加工）。

\end{enumerate}

\exercise{连接想法}{两人一组}{5}

这个练习是用精细加工改善保持的一个例子。
找一个同伴，
每人各自独立选一个想法，
然后亮出你们的想法，试着找出一条四环的链条，
从一个通向另一个。
例如，
如果两个想法是「萨斯喀彻温」和「统计」，
这些链环可能是：

\begin{itemize}

\item
  萨斯喀彻温是加拿大的一个省；

\item
  加拿大是一个国家；

\item
  国家有政府；

\item
  政府用统计来分析民意。

\end{itemize}

\exercise{趋同演化}{两人一组}{15}

上面没讲到的一种做法是\gref{g:guided-notes}{引导式笔记}，
它是老师准备好的笔记，
用来提示学习者在讲座或讨论中对关键信息作出回应。
提示可以是让学习者填补的空格、
紧跟在学习者应下定义术语旁的星号，
等等。

为你最近教过或打算教的一节课，做两到四张引导式笔记卡。
和同伴交换卡片：
看懂它要求做什么容易吗？
填完这些提示要花多长时间？
用在编程例子上效果如何？

\exercise{改变想法}{两人一组}{10}

\cite{Kirs2013}主张，关于「数字原住民」、
学习风格、
和自我教育者的种种迷思，全都反映了同一种错误信念，
即学习者知道自己最需要什么；
他还警告说，我们可能正处在一个向下的螺旋里，
教育研究者每一次试图驳斥这些迷思，
都反而让对手更加相信学习科学是伪科学。
从你到目前为止在本书中学到的关于学习的东西里，
挑一件让你吃惊、或与你此前相信的某事相矛盾的事，
练习用一两分钟把它讲给同伴听。
你讲得有说服力吗？

\exercise{卡片}{个人}{15}

用便利贴或手边任何别的东西，
为你最近教过或学过的某个主题做六七张卡片。
和同伴交换，看看你们各自要多久
才能达到 100\% 的完美回忆率。
做完后把卡片搁到一边，
半小时后再回来，看看你的回忆率是多少。

\exercise{使用 ADEPT}{全班}{15}

挑一件你最近教过或被教过的东西，
用 ADEPT 方法的五个步骤，简要设计一节介绍它的课。

\exercise{多任务的代价}{两人一组}{10}

\hreffoot{http://www.learningscientists.org/blog/2017/7/28-1}{Learning Scientists 博客}
描述了一个只用一块秒表就能做的简单实验，
用来展示多任务的心理代价。
两人一组，
测量每个人做下面这三项任务各需要多长时间：

\begin{itemize}
\item
  从 1 数到 26，数两遍。
\item
  从 A 到 Z 背诵字母表，背两遍。
\item
  把数字和字母交错起来，
  也就是说，「1、A、2、B、{\ldots}」
  这样下去。
\end{itemize}

让每对报告他们的数据。
在没有针对性练习的情况下，
第三项任务总是比其中任何一项单独任务都要显著更长。

\exercise{计算机教育中的迷思}{全班}{20}

\cite{Guzd2015b}列出了关于计算机教育的十大错误信念，
其中包括：

\begin{enumerate}
\item
  计算机科学里女性偏少，就跟其他所有 STEM 领域一样。
\item
  要让更多女性进入计算机科学，我们需要更多女性计算机科学教师。
\item
  学生评价是评估教学的最好办法。
\item
  好老师会按学生的学习风格来个性化教学。
\item
  好的计算机科学老师应该以身作则展示良好的软件开发实践，
  因为他们的职责是培养优秀的软件工程师。
\item
  有些人天生就比别人更擅长编程。
\end{enumerate}

让每个人对每一条投票 +1（同意）、-1（不同意）或 0（不确定），
然后阅读
\hreffoot{https://cacm.acm.org/blogs/blog-cacm/189498-top-10-myths-about-teaching-computer-science/fulltext}{原文}
里的完整解释，再投一次票。
哪些条目上人们的想法改变了？
哪些他们仍然认为是对的？为什么？

\exercise{校准式同伴互评}{两人一组}{20}

\begin{enumerate}

\item
  创建一份 5—10 分的评分标准，
  条目诸如「变量命名良好」「没有冗余代码」「控制流嵌套得当」，
  用来给希望学习者写出的那类程序打分。

\item
  挑选或创建一个包含 3—4 处违反上述条目的小程序。

\item
  按你的评分标准给这个程序打分。

\item
  让同伴用同一份评分标准给同一个程序打分。
  他接受了哪些你没有接受的？
  他批评了哪些你没有批评的？

\end{enumerate}

\section*{回顾}

\figpdfhere{figures/conceptmap-active-learning.pdf}{概念：主动学习}{f:individual-concept-map}{}
